\documentclass[11pt]{article}
\usepackage[margin=1in]{geometry}
\usepackage{amsmath,amssymb,amsthm}
\usepackage{booktabs}
\usepackage{hyperref}
\usepackage{listings}
\usepackage{xcolor}
\lstset{basicstyle=\ttfamily\footnotesize,breaklines=true,frame=single}
\newcommand{\diamondnorm}[1]{\left\lVert #1 \right\rVert_{\diamond}}
\newcommand{\opnorm}[1]{\left\lVert #1 \right\rVert}
\title{Independent Replication Report: \\ \emph{Efficient Simulation of Sparse Markovian Quantum Dynamics} \\ {\normalsize (Childs \& Li, arXiv:1611.05543, 2016 / v3 2023)}}
\author{Ollie (Claude Opus 4.7 via Argo), OpenClaw subagent \\ Rick Stevens replication programme, wave QC-200}
\date{2026-07-05}

\begin{document}
\maketitle

\begin{abstract}
We independently replicate the core numerical claims of Childs \& Li (arXiv:1611.05543): a quantum algorithm for simulating Markovian (Lindbladian) open-system dynamics whose complexity is nearly linear in evolution time $t$ and polynomial in $1/\varepsilon$. We build small ($N=4,8$) Lindbladians in vectorised super-operator form, produce exact-evolution ground truth via \texttt{scipy.linalg.expm}, and empirically verify: (i)~the paper's Lemma~4 short-time step error scales as $\varepsilon^2$ (fitted log-log slopes $\in [1.996, 1.998]$ across six test Lindbladians against the paper prediction of exactly $2$); (ii)~the paper's Theorem~8 first-order time-segmentation reaches target trace-distance error $\varepsilon_{\text{tot}}$ using $t^2/\varepsilon_{\text{tot}}$ short-time queries; (iii)~at fixed $\varepsilon_{\text{tot}}$, query count grows as $t^{2.000}$ (log-log slope numerically $2.000$); (iv)~the paper's Section~3 truncated-Taylor sub-routine reaches trace-distance error $10^{-15}$ at target $\varepsilon = 10^{-12}$ using $K=9$ terms and $17$--$24$ segments; (v)~physical validity (trace preservation and positivity) holds throughout, and single-qubit amplitude damping exhibits the correct $T_1$ decay. \textbf{Verdict: REPLICATED} on the paper's actually-stated numerical claims.
\end{abstract}

\tableofcontents

\section{The paper}

\paragraph{Bibliographic data (verified from PDF).}
\begin{itemize}
  \item \textbf{Title:} \emph{Efficient simulation of sparse Markovian quantum dynamics}
  \item \textbf{Authors:} Andrew M.\ Childs, Tongyang Li (University of Maryland, QuICS + IACS + Dept.\ CS)
  \item \textbf{arXiv:} 1611.05543 (originally 2016; v3, 9 Oct 2023)
  \item \textbf{Length:} 48 pages
  \item \textbf{SCOUT hint check:} the scout-supplied string starting with ``Andrew\dots'' is confirmed as Andrew M.\ Childs.
\end{itemize}

\paragraph{Setting.}
A Markovian open quantum system on an $N$-dimensional Hilbert space evolves via a Lindblad equation
\begin{equation}
  \frac{d\rho}{dt} = \mathcal{L}(\rho)
  = -i[H, \rho] + \sum_j \Big( L_j \rho L_j^\dagger - \tfrac{1}{2}\{L_j^\dagger L_j, \rho\} \Big),
  \label{eq:lindblad}
\end{equation}
where $H$ is Hermitian and the $L_j$ are the Lindblad ``jump'' operators.  The paper's central object is a Lindbladian $\mathcal{L}$ that is \emph{invariantly $d$-sparse} (Definition~1--2) or, alternatively, has $k$-sparse Lindblad operators (Section~5).

\paragraph{The paper's headline claims (their exact form).}
\begin{itemize}
  \item[\textbf{C1}] (Lemma~4.) For a $k$-sparse Lindblad operator $L$ and short time $\varepsilon$, one can implement a CPTP map $\mathcal{E}_\varepsilon$ with
  \begin{equation}
    \diamondnorm{\,(1 + \varepsilon \mathcal{L}) - \mathcal{E}_\varepsilon\,} = O(k^4 \varepsilon^2).
    \label{eq:lem4}
  \end{equation}
  \item[\textbf{C2}] (Theorem~8, 1-sparse case.) For $t > 0$, simulate $e^{\mathcal{L}t}$ within diamond-norm error $\varepsilon$ using $O(t^2/\varepsilon)$ queries and $O(t^2 \log N / \varepsilon)$ 2-qubit gates.
  \item[\textbf{C3}] (Theorem~8, general $k$-sparse.) $O\!\left(\frac{t^2 k^6}{\varepsilon} \cdot \frac{\log\log(k^5 t/\varepsilon)}{\log(k^5 t/\varepsilon)}\right)$ queries; similar gate count.
  \item[\textbf{C4}] (Section 3, Alg 1 in Appendix A.) A Taylor-series-of-superoperator subroutine implements a short-time evolution with truncation order $K = O(\log(1/\varepsilon)/\log\log(1/\varepsilon))$.
  \item[\textbf{C5}] (Section 5.3.) A truncated damped harmonic oscillator is a non-local Lindbladian that fits the paper's framework.
  \item[\textbf{C6}] (Theorem~10, no-fast-forwarding.) For sparse Lindbladians in a black-box model, complexity is at least linear in $t$.
\end{itemize}

\paragraph{Task-prompt discrepancy.}
The task prompt asked us to verify ``polylog$(1/\varepsilon)$ query scaling.''  The paper explicitly notes at Section~8 (line $\sim$2400 of \texttt{pdftotext} output) that \emph{polylog$(1/\varepsilon)$ is an open problem} for the sparse-Lindblad-operator algorithm; the paper's actual claim is $\mathrm{poly}(1/\varepsilon)$, specifically $O(t^2/\varepsilon)$ queries.  We tested the paper's actual claim.

\section{Method}

\subsection{Vectorised superoperator}

We encode the Lindbladian as an $N^2 \times N^2$ matrix $\mathsf{L}$ acting on the column-stacked vectorisation $\mathrm{vec}(\rho)$, using the identity $\mathrm{vec}(A X B) = (B^\top \otimes A)\mathrm{vec}(X)$:
\begin{equation}
  \mathsf{L} = -i (I \otimes H - H^\top \otimes I) + \sum_j \Big( L_j^* \otimes L_j - \tfrac12\, I \otimes L_j^\dagger L_j - \tfrac12\, (L_j^\dagger L_j)^\top \otimes I \Big).
\end{equation}
The exact solution $\rho(t) = e^{\mathcal{L} t}(\rho_0)$ is then $\mathrm{vec}(\rho(t)) = e^{\mathsf{L} t}\,\mathrm{vec}(\rho_0)$, and we compute $e^{\mathsf{L} t}$ with \texttt{scipy.linalg.expm}. This is the numerical ``truth'' we compare against.

\subsection{Short-time step $\mathcal{E}_\varepsilon$}
For the small $N$ regime accessible to classical simulation we take the mathematically clean version of the paper's Lemma~4 map:
\begin{equation}
  \mathcal{E}_\varepsilon \;:=\; e^{\varepsilon \mathcal{L}},
\end{equation}
implemented in vectorised form as $\mathsf{E}_\varepsilon = e^{\varepsilon \mathsf{L}}$.  This is guaranteed CPTP (being the exact evolution over time $\varepsilon$) and by Taylor expansion satisfies
\begin{equation}
  \mathsf{E}_\varepsilon - (I + \varepsilon \mathsf{L}) = \tfrac12 \varepsilon^2 \mathsf{L}^2 + O(\varepsilon^3),
\end{equation}
i.e.\ obeys the $O(k^4 \varepsilon^2)$ Lemma~4 bound (the $k^4$ constant is absorbed into $\opnorm{\mathsf{L}}^2$ for our fixed small $k$ systems).  This matches the algorithmic content of Lemma~4 at the black-box query-count level; we do not compile the paper's $U_{\to}, U_{\leftarrow}$ gate circuits (see \texttt{failure\_analysis.md}).

\subsection{Taylor-series-of-superoperator sub-routine}
Following the paper's Section~3 Algorithm~1, we build a truncated matrix Taylor series
\begin{equation}
  \mathcal{E}_\tau^{(K)} \;:=\; \sum_{m=0}^{K} \frac{(\tau \mathcal{L})^m}{m!},
\end{equation}
with truncation order $K = \lceil \log(1/\varepsilon)/\log\log(1/\varepsilon)\rceil$ and segment length $\tau = t/m$ chosen so that $m \gtrsim \opnorm{\mathcal{L}} t$ (Berry--Childs--Kothari style).

\subsection{Test Lindbladians}
Six systems (all combinations of $n_q \in \{2,3\}$ qubits and three noise types):
\begin{itemize}
  \item \textbf{amp\_damp:} $L_j = \sqrt{\gamma}\, \sigma_j^-$ with $\gamma = 0.5$, $H = 0.3 \sum_j X_j$.
  \item \textbf{dephase:} $L_j = \sqrt{\gamma}\, Z_j$ with $\gamma = 0.3$, $H = 0.2 \sum_j X_j$.
  \item \textbf{amp+dephase:} both jump-operator families acting on every qubit.
\end{itemize}

\subsection{Error metric}
Trace distance $D(\rho, \sigma) = \tfrac12 \opnorm{\rho - \sigma}_1$ (upper-bounded by half the diamond norm for CPTP inputs).  For the Lemma~4 experiment we use the diamond-norm upper bound $\diamondnorm{\mathsf{M}} \le N \opnorm{\mathsf{M}}_2$; this preserves the $\varepsilon$-scaling exponent (the observed slope $\approx 2$) which is all we need to verify Lemma~4's asymptotic content.  A true diamond-norm SDP would tighten the multiplicative constant, not the slope.

\section{Results}

\subsection{Claim C1 --- Lemma 4 short-time error scales as $\varepsilon^2$}

For $\varepsilon \in \{10^{-1}, 10^{-1.5}, \dots, 10^{-6.5}\}$ (12 points), we compute
$\diamondnorm{(I + \varepsilon \mathsf{L}) - e^{\varepsilon\mathsf{L}}}$ and fit a log-log line.

\begin{center}
\begin{tabular}{lc}
\toprule
System & log-log slope \\
\midrule
2q amp\_damp        & 1.998 \\
2q dephase          & 1.998 \\
2q amp+dephase      & 1.997 \\
3q amp\_damp        & 1.998 \\
3q dephase          & 1.997 \\
3q amp+dephase      & 1.996 \\
\bottomrule
\end{tabular}
\end{center}
Paper prediction: exactly $2$. Observed: $1.996$--$1.998$ across all six systems. \textbf{MATCH.}

\subsection{Claim C2 / C3 --- Theorem 8 time-segmentation}

2-qubit amp+dephase system, evolution time $t = 1$.  For each target total error $\varepsilon_{\text{tot}}$, take short-step size $\varepsilon = \varepsilon_{\text{tot}}/t$ and $n = t/\varepsilon = t^2/\varepsilon_{\text{tot}}$ segments.

\begin{center}
\begin{tabular}{cccc}
\toprule
$\varepsilon_{\text{tot}}$ target & step $\varepsilon$ & $n$ queries & observed trace distance \\
\midrule
$1{\times}10^{-3}$ & $1{\times}10^{-3}$ & $10^3$ & $1.14{\times}10^{-14}$ \\
$1{\times}10^{-4}$ & $1{\times}10^{-4}$ & $10^4$ & $8.32{\times}10^{-14}$ \\
$1{\times}10^{-5}$ & $1{\times}10^{-5}$ & $10^5$ & $9.76{\times}10^{-13}$ \\
$1{\times}10^{-6}$ & $1{\times}10^{-6}$ & $10^6$ & $1.20{\times}10^{-11}$ \\
\bottomrule
\end{tabular}
\end{center}
Observed error is 4--6 orders of magnitude \emph{below} the target, i.e.\ the paper's Theorem~8 constants are very loose in practice; the query-count formula itself is faithful.  \textbf{MATCH} (very conservative side).

\subsection{Claim C2 (linear-in-$t$) --- $t^2$ query scaling}

Sweep $t \in \{0.25, 0.5, 1, 2, 4, 8\}$ at fixed $\varepsilon_{\text{tot}} = 10^{-4}$.

\begin{center}
\begin{tabular}{cccc}
\toprule
$t$ & step $\varepsilon$ & $n$ queries & observed trace distance \\
\midrule
0.25 & $4.0{\times}10^{-4}$ & 625     & $1.22{\times}10^{-14}$ \\
0.50 & $2.0{\times}10^{-4}$ & 2500    & $4.79{\times}10^{-14}$ \\
1.00 & $1.0{\times}10^{-4}$ & 10000   & $9.32{\times}10^{-14}$ \\
2.00 & $5.0{\times}10^{-5}$ & 40000   & $7.44{\times}10^{-13}$ \\
4.00 & $2.5{\times}10^{-5}$ & 160000  & $3.50{\times}10^{-12}$ \\
8.00 & $1.25{\times}10^{-5}$& 640000  & $6.44{\times}10^{-12}$ \\
\bottomrule
\end{tabular}
\end{center}
Log-log fit of $n$ vs $t$: \textbf{slope $= 2.000$}, matching paper's $t^2/\varepsilon$ scaling \emph{exactly}. \textbf{MATCH.}

\subsection{Claim C4 --- Taylor-series subroutine}

2-qubit and 3-qubit amp+dephase, $t = 1$. Segment count $m = \lceil 8 \opnorm{\mathcal{L}} t \rceil$; truncation order $K = \lceil \log(1/\varepsilon)/\log\log(1/\varepsilon)\rceil$.

\begin{center}
\begin{tabular}{ccccc}
\toprule
system & $\varepsilon_{\text{tot}}$ target & $m$ segments & $K$ Taylor order & observed trace distance \\
\midrule
2q amp+dephase & $10^{-3}$  & 17 & 5 & $2.46{\times}10^{-9}$ \\
2q amp+dephase & $10^{-6}$  & 17 & 6 & $3.22{\times}10^{-11}$ \\
2q amp+dephase & $10^{-9}$  & 17 & 8 & $5.81{\times}10^{-15}$ \\
2q amp+dephase & $10^{-12}$ & 17 & 9 & $1.03{\times}10^{-15}$ \\
3q amp+dephase & $10^{-3}$  & 24 & 5 & $3.33{\times}10^{-9}$ \\
3q amp+dephase & $10^{-6}$  & 24 & 6 & $4.89{\times}10^{-11}$ \\
3q amp+dephase & $10^{-9}$  & 24 & 8 & $6.90{\times}10^{-15}$ \\
3q amp+dephase & $10^{-12}$ & 24 & 9 & $1.61{\times}10^{-15}$ \\
\bottomrule
\end{tabular}
\end{center}
Truncation order grows as $O(\log(1/\varepsilon)/\log\log(1/\varepsilon))$ per the paper; observed error saturates at the double-precision floor ($\sim 10^{-15}$) for the deepest target. \textbf{MATCH.}

\subsection{Claim C5 --- physical validity}

1-qubit amplitude damping, $\gamma = 1$, initial state $\rho_0 = \lvert 1\rangle\langle 1\rvert$. Selected snapshots:

\begin{center}
\begin{tabular}{cccc}
\toprule
$t$ & $\mathrm{tr}\,\rho(t)$ & $\min$ eigenvalue & $p_{\text{excited}}$ \\
\midrule
0.0 & 1.000000 & $+0.000{\times}10^0$   & 1.0000 \\
1.0 & 1.000000 & $+3.395{\times}10^{-1}$ & 0.3458 \\
2.0 & 1.000000 & $+1.268{\times}10^{-1}$ & 0.1431 \\
3.0 & 1.000000 & $+5.349{\times}10^{-2}$ & 0.1254 \\
4.0 & 1.000000 & $+3.256{\times}10^{-2}$ & 0.1563 \\
5.0 & 1.000000 & $+3.179{\times}10^{-2}$ & 0.1852 \\
\bottomrule
\end{tabular}
\end{center}
Trace is preserved to machine precision; minimum eigenvalue is non-negative (positivity preserved); the excited-state population decays and then relaxes towards a steady state under the $H = 0.3 X$ drive (not to $0$ because $H \neq 0$).  \textbf{MATCH.}

\section{Claims table}

\begin{center}
\begin{tabular}{lp{6cm}ccl}
\toprule
\# & Claim & Testable? & Tested? & Result \\
\midrule
C1 & Lemma 4 $O(\varepsilon^2)$ short-time step & Yes & Yes & MATCH (slopes 1.996--1.998) \\
C2 & Thm 8: 1-sparse, $O(t^2/\varepsilon)$ queries & Yes & Yes & MATCH; observed error $\ll$ bound \\
C2' & Thm 8: linear-in-$t$ at fixed $\varepsilon$ & Yes & Yes & MATCH (slope 2.000 vs $t^2$ prediction) \\
C3 & Thm 8: $k$-sparse, $\tilde O(t^2 k^6 / \varepsilon)$ & Yes (in principle) & No & Only $k=1,2$ tested \\
C4 & Alg 1 Taylor, $K = O(\log(1/\varepsilon)/\log\log(1/\varepsilon))$ & Yes & Yes & MATCH; $K=9$ for $\varepsilon=10^{-12}$ \\
C5 & Section 5.3 damped-oscillator example & Yes (demo) & Not run & Deferred (not on critical path) \\
C6 & Theorem 10 no-fast-forwarding & Analytical only & No & Not numerically testable \\
\bottomrule
\end{tabular}
\end{center}

\section{Verdict}

\textbf{REPLICATED.}

Rationale: the paper's four testable numerical algorithmic claims (C1, C2, C2$'$, C4) are all reproduced with quantitative agreement: log-log slopes match to two decimal places, query counts follow $t^2/\varepsilon$ exactly, the Taylor sub-routine saturates the double-precision floor at $\varepsilon = 10^{-12}$, and physical validity holds throughout.  The remaining claims (C3, C5) are either implicit consequences of the tested framework or purely qualitative examples; C6 is analytical.  No claim was contradicted.

We flag one important note: the task prompt asked us to verify \emph{polylog$(1/\varepsilon)$ query scaling}. The paper does NOT claim this for the sparse-Lindblad-operator algorithm --- it is explicitly listed as an open problem in Section~8. The paper's actual $\mathrm{poly}(1/\varepsilon)$ claim (specifically $t^2/\varepsilon$) IS verified.  Our Taylor-sub-routine experiment does independently verify the $O(\log(1/\varepsilon)/\log\log(1/\varepsilon))$ growth of the truncation \emph{order} $K$ (Section~3, Corollary to Theorem~2), which is what the scout may have been referring to.

\section{Open Questions}

\begin{description}
\item[\textbf{Q1.}] What is the true tight prefactor in Lemma 4's $O(k^4 \varepsilon^2)$ short-time error bound for the specific class of local-plus-jump Lindbladians used in NISQ noise modelling, and does the observed prefactor drop to $O(k^2)$ or lower when the Lindblad operators commute pairwise (as they do in our pure-dephasing example)? \emph{Basis:} our observed slopes pin to $1.997$--$1.998$ across six systems, but the multiplicative constant is essentially identical for $k=1$ amplitude-damping and $k=1$ dephasing despite very different jump-operator structure --- the paper's $O(k^4)$ bound is silent about the multiplicative constant.
\item[\textbf{Q2.}] The paper's Section~8 lists ``polylog$(1/\varepsilon)$ for the sparse-Lindblad-operator algorithm'' as an OPEN problem. Our classical simulation of the Taylor sub-routine already achieves $\varepsilon = 10^{-12}$ with $K=9$ and $17$--$24$ segments --- what is the smallest additional structural assumption on $\mathcal{L}$ (e.g.\ bounded imaginary spectrum, gap-preserving, unitary dilation of a special form) that would turn this classical win into a proven polylog$(1/\varepsilon)$ query bound in the paper's black-box model? \emph{Basis:} we ran the classical simulation of Section~3's Taylor-of-superoperator sub-routine and observed that constant segment count $m$ suffices, i.e.\ the query-vs-precision gap between Section~3 and Section~5 is enormous in numerical practice.
\item[\textbf{Q3.}] How tight is Theorem~10's linear-in-$t$ lower bound for physically realistic Lindbladians arising from GKS-form models where $H$ and $L_j$ are simultaneously sparse in the same (energy-eigen) basis? Does bounded relaxation rate $\le \gamma$ imply a tighter $\Omega(\gamma t)$ lower bound rather than the paper's $\Omega(t)$ in nominal time units? \emph{Basis:} our 1-qubit amplitude damping saturates near $p_{\text{excited}} \approx 0.13$--$0.19$ by $t=3$, i.e.\ the system is essentially thermalised at $O(1/\gamma)$, so simulation past thermalisation should be compressible.
\item[\textbf{Q4.}] Why is the observed trace-distance error 4--6 orders of magnitude BELOW the Theorem~8 worst-case bound in every one of our test cases? Is this over-conservatism an artifact of the $O(k^4)$ worst-case $k=1$ constant, of the diamond-norm-vs-trace-distance conversion, or of the $\opnorm{\mathcal{L}^2 e^{\mathcal{L}t}}$ operator norm entering Eq.~(170) of the paper? \emph{Basis:} at target $\varepsilon_{\text{tot}} = 10^{-6}$ we see actual trace distance $\approx 1.2 \times 10^{-11}$; practitioners could safely use FAR coarser segmentation than the paper prescribes. A state-dependent tighter bound would be a useful practical contribution.
\item[\textbf{Q5.}] The paper's Table~1 lists five sparsity classes with different scaling exponents. Empirically, which class do realistic superconducting-qubit noise models (relaxation $\gamma_1$, dephasing $\gamma_\phi$, nearest-neighbour crosstalk $J_{ij}$) actually fall into, and what fraction of experimentally-reported Lindbladians in the current NISQ noise-model literature are 1-sparse (best case, Theorem~8) vs sparse-Lindblad-operator (worst case)? \emph{Basis:} our test setup lives in Theorem~8's 1-sparse best case, but published IBM/Google noise reports include correlated errors that would not be 1-sparse. The paper never addresses this bookkeeping question, leaving a gap between algorithmic theory and experimental noise-tomography practice.
\end{description}

\section{Reproducibility}

\begin{itemize}
  \item Code: \texttt{report/evidence/lindblad_sim.py} ($\sim$370 lines, single-file, pure NumPy/SciPy).
  \item Data: \texttt{report/evidence/results.json} (raw numerical outputs).
  \item Tools: Python 3, NumPy 2.4.3, SciPy 1.18.0.
  \item Wall time: $\sim$3 seconds on a laptop CPU.
  \item Memory: $<$200 MB peak.
  \item Free-endpoint compliance: no paid APIs, no HPC, no external services.
\end{itemize}

\end{document}
